% TEMPLATE-MILC-v3.tex - plantilla maestra UNICA de las guias MILC v3.
%
% La usan las cuatro familias de guias, cada una con su driver:
%   · Grados 6-11          -> scripts/build-guias-g11.py
%   · Bebras               -> scripts/build-guias-bebras.py
%   · Semillero            -> scripts/build-guias-semillero.py
%   · Territorio interior  -> scripts/build-guias-territorio-interior.py
%
% Antes habia tres copias casi identicas (~400 lineas iguales, 6 distintas):
% cada mejora habia que aplicarla tres veces, y en la practica no se hacia
% (el motor de recursos vivia solo en dos de ellas). Lo que varia entre
% familias esta parametrizado en 6 placeholders que cada driver llena
% (se nombran aqui SIN su sintaxis real: el builder reemplaza por texto plano,
% tambien dentro de los comentarios, y un valor multilinea rompe el preambulo):
%
%   HEADER_IZQ        encabezado izquierdo de pagina
%   HEADER_DER        encabezado derecho de pagina
%   PORTADA_ETIQUETA  rotulo sobre el numero grande (GRADO/MOMENTO/MODULO)
%   PORTADA_SERIE     linea de serie de la portada
%   PORTADA_ID        identificador de la guia en la portada
%   PORTADA_CIRCULO   el circulito de grado (°), o vacio si no aplica
%   PDF_TITULO        titulo en texto plano (sin \textbf ni comillas TeX) para
%                     los metadatos del PDF (/Title); lo produce lib_xelatex.texto_pdf
%
% Composicion (auditoria 2026-09-03): las bandas de estacion piden espacio
% (\Needspace*) y prohiben el salto justo despues (\nopagebreak), asi no quedan
% huerfanas al pie; las cajas largas (softbox, drawbox, infoband, puentebox) son
% `breakable`, asi una caja de media pagina ya no arrastra todo a la siguiente
% dejando la anterior al 40 %. Todo texto sobre mostaza o verde va en milcNegro
% (blanco sobre #E5B400 da 1,93:1 y sobre #58A923 2,95:1; ninguno pasa WCAG AA).
% El resto de claves las documenta content/guias/_SCHEMA.md. El script valida
% que no queden placeholders sin reemplazar antes de escribir el archivo final.

% Etiquetado PDF (latex-lab): /Lang, estructura y texto alternativo de las
% figuras, para lectores de pantalla. Debe ir ANTES de \documentclass.
\DocumentMetadata{lang=es-CO,tagging=on}
\documentclass[11pt,letterpaper]{article}
% headheight=24pt: la cabecera derecha lleva dos lineas (identidad de la guia
% y estacion actual). headsep se acorta en la misma medida para que el bloque
% de texto quede exactamente donde estaba con headheight=12pt.
\usepackage[margin=2.0cm,headheight=24pt,headsep=13pt]{geometry}
\usepackage[table]{xcolor}
\usepackage{fontspec}
\usepackage[spanish]{babel}
\usepackage{tikz}
% Librerías TikZ para los diagramas de las guías (bloque `recursos:`):
%  · babel  → babel en español vuelve activos caracteres como «>», y TikZ los usa
%             en sus opciones (>=stealth, ->). Sin esta librería, cualquier
%             diagrama con flechas revienta con «\language@active@arg> has an
%             extra }». Debe ir ANTES de que se dibuje nada.
%  · positioning        → sintaxis «below left=2pt and 3pt».
%  · decorations.pathreplacing → llaves ({brace}) para señalar rangos.
\usetikzlibrary{babel,positioning,decorations.pathreplacing}
\usepackage{graphicx}
% Circuitos: el trabajo de electrónica se hace hoy con micro:bit y sus sensores
% integrados (MakeCode, sin cableado), así que no se necesitan esquemáticos.
% Si algún día entran montajes con protoboard, basta descomentar esta línea:
% los diagramas .tex/.tikz declarados en `recursos:` ya se incrustan solos.
% \usepackage{circuitikz}
\usepackage[most]{tcolorbox}
\usepackage{tabularx}
\usepackage{array}
\usepackage{fancyhdr}
\usepackage{titlesec}
\usepackage[document]{ragged2e}  % bandera a la derecha en todo el cuerpo
\usepackage{enumitem}
\usepackage{microtype}
\usepackage{etoolbox}
\usepackage{needspace}
% hyperref al final del preambulo: metadatos del PDF (titulo, autor, idioma,
% familia), marcadores por estacion (\pdfbookmark) y enlaces sin recuadro.
\usepackage[hidelinks,
  pdftitle={El aula que excluye: barreras que nadie puso a propósito},
  pdfauthor={PhD. Álvaro Cárdenas Orozco},
  pdfsubject={Territorio interior · Educación socioemocional · Grado 9° · Momento 6 · MILC v3},
  pdflang={es-CO},
  bookmarksopen=true,bookmarksnumbered=false]{hyperref}

% Helvetica Neue no trae la flecha U+2192, asi que cada «→» escrito literalmente
% en el YAML se compilaba a NADA, en silencio (462 flechas en 61 guias: rutas del
% tipo «paleta → tipografia → lockup» salian rotas). Mapeamos el caracter a su
% version matematica, que si tiene glifo. Asi el autor puede seguir escribiendo
% «→» comodamente en el YAML.
\usepackage{newunicodechar}
\newunicodechar{→}{\ensuremath{\rightarrow}}
% Mismo problema con los simbolos de marca y vinetas: el «✓» de «una palomita ✓»
% salia como cuadrito vacio en el PDF de todas las guias que lo usaban.
\usepackage{pifont}
\newunicodechar{✓}{\ding{51}}
\newunicodechar{✗}{\ding{55}}
\newunicodechar{✔}{\ding{52}}
\newunicodechar{•}{\textbullet}
\newunicodechar{≥}{\ensuremath{\geq}}
\newunicodechar{≤}{\ensuremath{\leq}}

\setmainfont{Helvetica Neue}
\setsansfont{Helvetica Neue}
\setlength{\parindent}{0pt}
\setlength{\parskip}{6pt}
% Sin particion de palabras: en 6.º-7.º una palabra cortada por guion al final
% de linea («Comuni-cacion») frena la lectura mucho mas que una linea corta.
\hyphenpenalty=10000
\exhyphenpenalty=10000
\pretolerance=2000
\tolerance=3000
\setlength{\emergencystretch}{3em}
\linespread{1.10}
\renewcommand{\arraystretch}{1.25}
\setlist{leftmargin=5mm,itemsep=1mm,topsep=1mm}
\newcolumntype{Y}{>{\RaggedRight\arraybackslash}X}

\definecolor{milcMagenta}{HTML}{E6007E}
\definecolor{milcVino}{HTML}{5A0038}
\definecolor{milcTurquesa}{HTML}{0093A5}
\definecolor{milcVerde}{HTML}{58A923}
\definecolor{milcMostaza}{HTML}{E5B400}
\definecolor{milcOcre}{HTML}{B86600}
\definecolor{milcNegro}{HTML}{241816}
\definecolor{milcGris}{HTML}{F7F7F4}
\definecolor{milcLinea}{HTML}{E8E2DD}
\definecolor{milcGradePrimary}{HTML}{0D9488}
\definecolor{milcGradeDark}{HTML}{115E59}
\definecolor{milcGradeSoft}{HTML}{CCFBF1}
\definecolor{milcGradeAccent}{HTML}{CFE7FF}
\definecolor{milcGradeText}{HTML}{FFFFFF}
\color{milcNegro}

\pagestyle{fancy}
\fancyhf{}
% Cabecera de dos lineas: identidad de la guia arriba y, a la derecha y en
% gris, la estacion en curso (\markright desde \milcestacion). Antes de la
% primera estacion la segunda linea va vacia; el \strut mantiene la altura
% para que la linea de arriba no baile entre paginas.
\lhead{\small Territorio interior · Educación socioemocional\\[-1pt]{\footnotesize\strut}}
\rhead{\small Grado 9° · Momento 6 · MILC v3\\[-1pt]{\footnotesize\strut\color{milcNegro!65}\rightmark}}
\cfoot{\small\thepage}
\renewcommand{\headrulewidth}{0pt}

% Sobre mostaza (#E5B400) y verde (#58A923) el texto va en milcNegro; sobre el
% resto de la paleta, en blanco. Se usa dentro de fonttitle/fontupper, donde un
% \color posterior manda sobre coltitle.
\newcommand{\milctintasobre}[1]{%
  \ifboolexpr{test{\ifstrequal{#1}{milcMostaza}} or test{\ifstrequal{#1}{milcVerde}}}%
    {\color{milcNegro}}{\color{white}}}

\newtcolorbox{titlebox}[1]{
  enhanced,colback=#1,colframe=#1,arc=7mm,boxrule=0pt,
  left=7mm,right=7mm,top=3mm,bottom=3mm,
  before skip=8pt,after skip=8pt,
  fontupper=\sffamily\bfseries\Large\color{white}
}

\newtcolorbox{softbox}[3]{
  enhanced,breakable,pad at break=3mm,
  colback=#2,colframe=#1,borderline west={4pt}{0pt}{#1},
  arc=2mm,boxrule=.4pt,left=4mm,right=4mm,top=3mm,bottom=3mm,
  before skip=6pt,after skip=8pt,title={#3},coltitle=white,
  colbacktitle=#1,fonttitle=\sffamily\bfseries\milctintasobre{#1},fontupper=\RaggedRight,
  attach boxed title to top left={xshift=3mm,yshift=-2mm},
  boxed title style={arc=2mm, boxrule=0pt}
}

\newtcolorbox{drawbox}[2]{
  enhanced,breakable,pad at break=4mm,
  colback=white,colframe=#1,arc=4mm,boxrule=1pt,
  left=5mm,right=5mm,top=4mm,bottom=4mm,before skip=8pt,after skip=8pt,
  title={#2},coltitle=white,colbacktitle=#1,fonttitle=\sffamily\bfseries\milctintasobre{#1},
  fontupper=\RaggedRight,
  attach boxed title to top left={xshift=4mm,yshift=-2mm},
  boxed title style={arc=3mm, boxrule=0pt}
}

\newtcolorbox{infoband}[1]{
  enhanced,breakable,pad at break=2mm,colback=#1!8,colframe=#1!50,arc=2mm,boxrule=.5pt,
  left=4mm,right=4mm,top=2mm,bottom=2mm,before skip=4pt,after skip=4pt,
  fontupper=\small\RaggedRight
}

% Puente narrativo entre secciones (contrato editorial Conéctate).
% Da continuidad: "Acabas de X. Ahora vas a Y porque Z."
% Visualmente: banda izquierda en color del grado, texto en itálica pequeña.
\newtcolorbox{puentebox}{
  enhanced,breakable,pad at break=2mm,colback=milcGradeSoft,colframe=milcGradeDark,
  borderline west={3pt}{0pt}{milcGradeDark},
  arc=0mm,boxrule=0pt,left=5mm,right=4mm,top=2.5mm,bottom=2.5mm,
  before skip=4pt,after skip=8pt,
  fontupper=\itshape\small\color{milcNegro}\RaggedRight
}
% La flecha va en modo matematico a proposito: Helvetica Neue NO tiene el glifo
% U+2192, asi que \textrightarrow se compilaba a NADA («Missing character: There
% is no → in font Helvetica Neue») y el puente salia sin flecha. En modo
% matematico la flecha viene de la fuente matematica y si se dibuja.
\newcommand{\puente}[1]{%
  \begin{puentebox}%
    \textcolor{milcGradeDark}{$\rightarrow$}\hspace{2mm}#1%
  \end{puentebox}%
}

\newcommand{\linead}[1]{\noindent\color{milcLinea}\rule{#1}{0.5pt}\color{milcNegro}}
\newcommand{\respuesta}[1]{\par\vspace{2mm}\foreach \n in {1,...,#1}{\noindent\color{milcLinea}\rule{\linewidth}{0.5pt}\par\vspace{5mm}}\color{milcNegro}}
\newcommand{\checkbox}{\textcolor{milcGradeDark}{\fbox{\phantom{x}}}\hspace{5pt}}

% ─── Listas aireadas para las guías socioemocionales ────────────────────
% Componer las verificaciones como «\checkbox texto\\» deja el texto que salta
% de línea debajo del cuadrito y apelmaza el bloque. Estos entornos alinean la
% continuación y airean los ítems. Los usa Territorio Interior; cualquier
% familia puede hacerlo.
\newlist{checklista}{itemize}{1}
\setlist[checklista]{
  label=\textcolor{milcGradeDark}{\fbox{\phantom{x}}},
  leftmargin=9mm, labelsep=3.2mm, itemsep=2.4mm,
  topsep=2.5mm, parsep=0pt, partopsep=0pt, before=\RaggedRight
}
\newlist{pasoslista}{enumerate}{1}
\setlist[pasoslista]{
  label=\textcolor{milcGradeDark}{\sffamily\bfseries\arabic*.},
  leftmargin=8mm, labelsep=3mm, itemsep=2.8mm,
  topsep=1mm, parsep=0pt, partopsep=0pt, before=\RaggedRight
}

\newcommand{\phasecard}[4]{
\begin{tcolorbox}[enhanced,colback=#3,colframe=#2,arc=3mm,boxrule=.5pt,height=3.05cm,valign=center,halign=center,left=1.5mm,right=1.5mm,top=2mm,bottom=2mm]
{\sffamily\bfseries\fontsize{22}{24}\selectfont\textcolor{#2}{#1}}\par
{\sffamily\bfseries\small #4}
\end{tcolorbox}
}

% ── Piezas del rediseno 2026 (legibilidad 6.º-11.º) ───────────────────
% Banda de estacion: reemplaza el titlebox macizo. Titulo a la izquierda,
% dato practico (tiempo/modalidad) a la derecha, en una sola linea.
% Nunca huerfana: pide 9 lineas libres (o salta de pagina rellenando con
% \vfill) y prohibe el salto justo despues de la banda. Nueve y no seis:
% la caja partible que sigue necesita ~80pt (titulo adosado, relleno y dos
% lineas) para arrancar en la misma pagina; con menos, tcolorbox la manda
% entera a la siguiente y la banda se queda sola. El penalty 10000 va
% en `after`, ANTES del espacio posterior: un `after skip` (aunque sea 0pt)
% es un \vskip tras la caja, y ese glue es un punto de corte legal que un
% \nopagebreak escrito despues de \end{tcolorbox} ya no cierra. Cada banda
% deja marcador PDF y actualiza la cabecera derecha.
\newcounter{milcestacion}
\newcommand{\milcestacion}[2]{%
\Needspace*{9\baselineskip}%
\stepcounter{milcestacion}%
\pdfbookmark[1]{#2}{milcestacion\themilcestacion}%
\markright{#2}%
\begin{tcolorbox}[enhanced,colback=#1,colframe=#1,arc=2mm,boxrule=0pt,
  left=5mm,right=5mm,top=2.6mm,bottom=2.6mm,before skip=11pt,
  after={\par\nopagebreak\vskip7pt}]
{\sffamily\bfseries\fontsize{15}{18}\selectfont\milctintasobre{#1}#2}
\end{tcolorbox}}

% Tarjeta de paso numerada: sustituye las tablas «Pilar / Aplicacion», que
% eran formato de consulta docente, no de estudiante. El numero va en un
% circulo sobre el borde para que la secuencia se lea de un vistazo.
\newcommand{\milcpaso}[3]{%
\Needspace{4\baselineskip}%
\begin{tcolorbox}[enhanced,colback=white,colframe=milcLinea,arc=1.5mm,boxrule=.9pt,
  left=14mm,right=4mm,top=3mm,bottom=3mm,before skip=4pt,after skip=4pt,
  fontupper=\RaggedRight,
  overlay={\node[anchor=north west,circle,fill=#2,text=white,inner sep=0pt,
    minimum size=7.4mm,font=\sffamily\bfseries\small]
    at ([xshift=3.6mm,yshift=-3.2mm]frame.north west) {#1};}]
#3
\end{tcolorbox}}

% Casilla marcable de tamano util con lapiz (la anterior era \fbox{\phantom{x}},
% de alto variable segun el texto que la seguia).
\newcommand{\milcmarca}{\framebox[3.8mm]{\rule{0pt}{3.4mm}}}

% Fila marcable de lista suelta (checks de la estacion 1).
\newcommand{\milccheckrow}[1]{%
\par\vspace{1.8mm}\noindent
\makebox[6.4mm][l]{\raisebox{-.55mm}{\textcolor{milcNegro}{\milcmarca}}}%
\parbox[t]{\dimexpr\linewidth-6.4mm\relax}{\RaggedRight #1}\par}

% Celda marcable dentro de la rubrica (tabla de autoevaluacion). Misma
% sangria francesa, pero pensada para vivir dentro de una columna X.
\newcommand{\milcopcion}[1]{%
\makebox[5.2mm][l]{\raisebox{-.55mm}{\textcolor{milcNegro}{\milcmarca}}}%
\parbox[t]{\dimexpr\linewidth-5.2mm\relax}{\RaggedRight #1}}

% ── Recursos visuales de la guía ──────────────────────────────────────
% Declarados en el bloque `recursos:` del YAML y emitidos por el builder.
% \guiaFigura   → imagen rasterizada o PDF (foto, carta celeste, montaje).
% \guiaDiagrama → fuente TikZ/circuitikz (.tex) incrustada como vector:
%                 es la vía para circuitos y esquemáticos editables.
% [#1] = texto alternativo (alt) · #2 = ruta relativa al .tex · #3 = pie
% El alt llega al PDF etiquetado (/Alt) y lo lee el lector de pantalla.
\newcommand{\guiaFigura}[3][]{%
\begin{center}
\includegraphics[width=0.88\linewidth,height=0.38\textheight,keepaspectratio,alt={#1}]{#2}\\[1.5mm]
{\footnotesize\itshape\textcolor{milcNegro!75}{#3}}
\end{center}
}

\newcommand{\guiaDiagrama}[2]{%
\begin{center}
\input{#1}\\[1.5mm]
{\footnotesize\itshape\textcolor{milcNegro!75}{#2}}
\end{center}
}

\begin{document}
\thispagestyle{empty}

% ── Portada ───────────────────────────────────────────────────────────
% Antes: un rectangulo de color a pagina completa. Se veia bien en pantalla,
% pero en la fotocopiadora del colegio se comia el toner y salia gris sucio.
% Ahora la banda ocupa el tercio superior y el resto va en blanco.
\begin{tikzpicture}[remember picture,overlay]
  \path[fill=milcGradePrimary]
    (current page.north west) rectangle ([yshift=-10.6cm]current page.north east);

  \node[anchor=north west,text=milcGradeText] at ([xshift=2.0cm,yshift=-1.55cm]current page.north west)
    {\sffamily\bfseries\fontsize{12}{14}\selectfont MOMENTO};
  \node[anchor=north west,text=milcGradeText,opacity=.85] at ([xshift=2.0cm,yshift=-2.35cm]current page.north west)
    {\sffamily\bfseries\fontsize{8.5}{10}\selectfont TERRITORIO INTERIOR · SOCIOEMOCIONAL};
  \node[anchor=north east,text=milcGradeText,opacity=.85,align=right] at ([xshift=-2.0cm,yshift=-1.55cm]current page.north east)
    {\sffamily\bfseries\fontsize{9}{12}\selectfont I.E. Sor María Juliana\\Cartago, Valle del Cauca};

  \node[anchor=west,text=milcGradeText] (gradonum) at ([xshift=1.85cm,yshift=-7.1cm]current page.north west)
    {\sffamily\bfseries\fontsize{112}{112}\selectfont 6};
  % El circulito de grado (°) solo aplica a las guias de grado («11°»); en
  % Bebras y Semillero el numero grande es un momento/modulo, no un grado.
  % Se posiciona relativo al nodo `gradonum`, no en coordenadas absolutas.
  

  \node[anchor=west,text=milcGradeText,text width=11.4cm,align=left] at ([xshift=1.5cm,yshift=0.5cm]gradonum.east)
    {
     {\sffamily\bfseries\fontsize{9}{11}\selectfont\MakeUppercase{Territorio interior · Socioemocional}}\\[3mm]
     {\sffamily\bfseries\fontsize{21}{25}\selectfont\setlength{\baselineskip}{27pt}El aula\\[4pt]que excluye\\[4pt]barreras que nadie puso}\\[3.5mm]
     {\sffamily\bfseries\fontsize{15}{18}\selectfont Grado 9° · Momento 6}
    };
\end{tikzpicture}

\vspace*{9.4cm}
\vfill

{\sffamily\bfseries\fontsize{8}{10}\selectfont\textcolor{milcGradeDark}{LO QUE TE LLEVAS HOY}}

\begin{tcolorbox}[enhanced,colback=white,colframe=milcNegro,arc=2mm,boxrule=1.6pt,
  left=5mm,right=5mm,top=4mm,bottom=4mm,before skip=3pt,after skip=12pt,
  fontupper=\RaggedRight\fontsize{11}{15.5}\selectfont]
Tu auditoría de barreras: el recorrido real de tu colegio con las barreras físicas, de comunicación y de método que encontraste, cada una con quién queda por fuera y con un ajuste que no cuesta plata, más la lista entregada a alguien que pueda hacer algo.
\end{tcolorbox}

\begin{tcolorbox}[enhanced,colback=milcGris,colframe=milcGris,arc=2mm,boxrule=0pt,
  left=5mm,right=5mm,top=3.5mm,bottom=3.5mm,after skip=12pt]
{\sffamily\bfseries\fontsize{8}{10}\selectfont\textcolor{milcNegro!55}{ESTA GUÍA ES DE}}\\[3mm]
\begin{tabularx}{\linewidth}{X p{3.2cm} p{3.2cm}}
\linead{\linewidth} & \linead{3.0cm} & \linead{3.0cm}\\[-1mm]
{\fontsize{8.5}{10}\selectfont\textcolor{milcNegro!55}{Nombre completo}} &
{\fontsize{8.5}{10}\selectfont\textcolor{milcNegro!55}{Grupo}} &
{\fontsize{8.5}{10}\selectfont\textcolor{milcNegro!55}{Fecha}}\\
\end{tabularx}
\end{tcolorbox}

\vfill

\noindent\textcolor{milcLinea}{\rule{\linewidth}{1pt}}\\[2mm]
{\sffamily\bfseries\fontsize{9.5}{12}\selectfont Autor: Dr. Álvaro Cárdenas Orozco}\\[1mm]
{\sffamily\fontsize{8.5}{11}\selectfont\textcolor{milcNegro!55}{Metodología MILC · Apertura ancestral · Discapacidad y barreras · Triángulo Dussel-Estoicismo-Floridi}}

\newpage

\section*{Apertura: El aula que excluye: barreras que nadie puso a propósito}

\begin{softbox}{milcGradeDark}{milcGradeSoft}{Saber ancestral}
En \textbf{1940}, en el barrio San Fernando de Cali, la caleña \textbf{Luisita Sánchez de Hurtado} fundó el \textbf{Instituto para Niños Ciegos y Sordos}. \emph{Fue una obra enorme y sigue siendo una de las instituciones más importantes del país en su campo.} \textbf{Y al mismo tiempo dice algo sobre la época que conviene leer con cuidado:} \emph{empezó como un \textbf{internado}}, al que asistían personas de todas las edades ---\textbf{es decir, la manera de atender a los niños ciegos y sordos era \emph{sacarlos} de la escuela común y reunirlos aparte}---. \emph{Con los años eso cambió:} tras los daños del terremoto de \textbf{1979} tuvo que cerrar y volvió convertido en \textbf{externado}. \textbf{Y falta el dato que completa el cuadro.} \emph{Los niños sordos de este país se comunicaban entre ellos desde siempre, con una lengua propia.} \textbf{Pero esa lengua no fue reconocida por el Estado colombiano como \emph{idioma} sino hasta la \textbf{Ley 324 de 1996}}, que reconoció la \textbf{Lengua Manual Colombiana} como lengua propia de la comunidad sorda y ordenó garantizar intérpretes y promover escuelas para formarlos. \emph{Fíjate en la fecha: 1996. Es decir, hubo generaciones enteras cuya lengua materna \textbf{oficialmente no era una lengua}.}
\end{softbox}

\begin{softbox}{milcTurquesa}{milcGris}{Saber contemporáneo}
¿Y dónde está exactamente la discapacidad: en la persona o en el lugar? \textbf{Hay una respuesta oficial, y Colombia la firmó.} La \textbf{Convención sobre los Derechos de las Personas con Discapacidad}, adoptada por la Asamblea General de las Naciones Unidas el \textbf{13 de diciembre de 2006} y aprobada en Colombia por la \textbf{Ley 1346 de 2009}, dice algo que cambia por completo dónde hay que mirar: \emph{«la discapacidad es un concepto que evoluciona y que \textbf{resulta de la interacción} entre las personas con deficiencias y \textbf{las barreras debidas a la actitud y al entorno} que evitan su participación plena y efectiva en la sociedad»}. \textbf{Lee otra vez esas dos palabras: \emph{interacción} y \emph{barreras}.} \emph{No dice que la discapacidad sea una característica de la persona.} \textbf{Dice que aparece \emph{entre} la persona y un entorno que fue construido sin contar con ella.} \emph{La consecuencia práctica es enorme y es el trabajo de hoy:} \textbf{si la barrera está en el entorno, entonces se puede quitar} ---y quitarla no es un favor que se le hace a nadie: es cumplir lo que ya está firmado---.
\end{softbox}

\begin{softbox}{milcMostaza}{milcGris}{Pregunta-puente}
\textit{Hubo generaciones cuya lengua materna \textbf{oficialmente no era una lengua}, y niños a los que se atendía \textbf{sacándolos} de la escuela común. \textbf{¿Qué hay en tu colegio ---en el edificio, en cómo se explica, en cómo se avisan las cosas--- que deja a alguien por fuera} \ldots \textbf{sin que nadie lo haya querido}?}
\end{softbox}

\begin{softbox}{milcVerde}{milcGris}{Saber y saber hacer}
Al terminar podrás: (1) \textbf{explicar} que la discapacidad resulta de la interacción entre una persona y las barreras del entorno, según la norma que Colombia firmó; (2) \textbf{distinguir} tres clases de barrera ---física, de comunicación y de método--- y reconocer que las dos últimas son las que casi nadie ve; (3) \textbf{auditar} tu propio colegio y encontrar barreras concretas; (4) \textbf{proponer} ajustes que no cuestan dinero y entregarlos a alguien que pueda decidir.
\end{softbox}



\small
\begin{tabularx}{\textwidth}{>{\bfseries}p{3.0cm}Y}
\rowcolor{milcGradePrimary!12}Autor & Dr. Álvaro Cárdenas Orozco\\
DBA & Comprende los fenómenos que afectan a su territorio, regula sus emociones ante ellos y acompaña a otros con empatía (transversal 6°--11°, articulado con la política «Escuela, territorio de vida» del MEN, Resolución 006519 de 2025, y la Ley 1620 de 2013)\\
Referentes & Primera ayuda psicológica (OMS, 2011/2012) · Directrices IASC (2007) · USGS y Servicio Geológico Colombiano · Pueblo nasa y la minga (CRIC) · Dussel · Estoicismo · Floridi\\
Producto final & Tu auditoría de barreras: el recorrido real de tu colegio con las barreras físicas, de comunicación y de método que encontraste, cada una con quién queda por fuera y con un ajuste que no cuesta plata, más la lista entregada a alguien que pueda hacer algo.\\
\end{tabularx}
\normalsize

\puente{En el momento 5 viste una jerarquía que se sostiene sin que nadie la defienda. \textbf{Hoy verás algo parecido y todavía más invisible:} \emph{barreras que excluyen sin que nadie las haya puesto con esa intención}.}

\milcestacion{milcGradeDark}{Ruta MILC: así vamos a aprender}
\begin{tabularx}{\textwidth}{>{\centering\arraybackslash}X>{\centering\arraybackslash}X>{\centering\arraybackslash}X>{\centering\arraybackslash}X}
\phasecard{1}{milcVerde}{milcGris}{Escucha\\[-1mm]\small Observo, escucho y pregunto.} &
\phasecard{2}{milcTurquesa}{milcGris}{Sistematización\\[-1mm]\small Miro el entorno, no a la persona.} &
\phasecard{3}{milcMagenta}{milcGris}{Praxis\\[-1mm]\small Construyo y aplico.} &
\phasecard{4}{milcVino}{milcGris}{Evaluación\\[-1mm]\small Reflexión liberadora.}
\end{tabularx}


\puente{Primero, dentro de tu propio salón: qué necesita alguien para poder participar, y qué le falta.}

\milcestacion{milcVerde}{Estación 1 de 3 · Escucha}

\begin{softbox}{milcVerde}{milcGris}{Escena para escuchar}
\textbf{Actividad 1 · IDENTIFICA --- Quién no puede hacer qué.} \textbf{Sin nombres: usa iniciales.} \textbf{Primera parte.} Piensa en tu salón y en tu colegio y anota \textbf{qué actividades le resultan difíciles o imposibles a alguien}: subir al segundo piso, oír lo que dice el profesor de espaldas, leer del tablero desde atrás, entender una explicación larga y hablada, escribir rápido, quedarse quieto una hora, entrar al baño, participar en la izada. \textbf{Segunda parte, la que cambia la mirada.} Al frente de cada una, escribe \textbf{qué hay en el lugar} que lo hace difícil ---no qué tiene la persona---. \emph{Es decir: en vez de «no puede subir», escribe «no hay ascensor ni rampa».} \textbf{Tercera parte.} ¿Hay alguien en tu colegio que \textbf{no participa de algo} y todo el mundo lo da por normal? \textbf{Y la última:} ¿alguna vez \textbf{tú} no entendiste algo y no preguntaste? \emph{¿Qué te lo impidió: no entender, o que preguntar costara?}
\end{softbox}

{\sffamily\bfseries\fontsize{8}{10}\selectfont\textcolor{milcNegro!55}{MARCA CUANDO LO HAYAS HECHO}}
\milccheckrow{Anotaste actividades concretas, no características de personas.}
\milccheckrow{Escribiste \textbf{qué hay en el lugar} que lo hace difícil, no qué tiene la persona.}
\milccheckrow{Identificaste a alguien que no participa de algo y todos lo dan por normal.}

\begin{infoband}{milcVerde}


\textbf{En tu cuaderno (Actividad 1 · IDENTIFICA).} Título: «Actividad 1. Quién no puede hacer qué». \textbf{Formato:} la lista de actividades con la barrera del entorno al frente + quién no participa + tu propia experiencia de no preguntar. \textbf{Extensión:} media página. \textbf{Sabes que terminaste cuando} todas tus frases hablan \textbf{del lugar} y no de las personas. Esta página es tuya: \textbf{nadie más tiene que leerla si tú no quieres}.
\end{infoband}

\puente{Ya lo tienes. Ahora, dónde está la discapacidad según la Convención, y por qué una lengua tardó hasta 1996 en ser reconocida.}

\milcestacion{milcTurquesa}{Fase 2 · Sistematización: entender la barrera}

\begin{softbox}{milcTurquesa}{milcGris}{Cuatro claves sobre las barreras}
Cuatro claves para dejar de mirar a la persona y aprender a mirar el lugar.
\end{softbox}

\milcpaso{1}{milcTurquesa}{\textbf{Aparte no es incluido, aunque se haga con cariño.} El Instituto de Cali nació en \textbf{1940} de una intención generosa \emph{y funcionaba como internado}: la forma de atender era \textbf{reunir aparte}. \emph{Y esto hay que decirlo con cuidado, porque no es un reproche a quien lo fundó ---en su época era lo que existía y era mejor que nada---.} \textbf{Pero muestra la lógica de una era entera:} \emph{si el niño no encaja en la escuela, se cambia al niño de escuela}. \textbf{Lo que cambió después es el orden de la frase:} \emph{si el niño no encaja en la escuela, \textbf{se cambia la escuela}}. \emph{Y ojo, porque esa lógica antigua sigue viva en versiones pequeñas:} \textbf{el que «hace las actividades aparte», el que se queda en el salón cuando los demás salen, el que tiene su propia hoja distinta y nunca la comparte con nadie.}}
\milcpaso{2}{milcTurquesa}{\textbf{La discapacidad está en la interacción, y eso lo dice la norma.} La Convención de \textbf{2006} ---aprobada aquí por la \textbf{Ley 1346 de 2009}--- define la discapacidad como algo que \textbf{resulta de la interacción} entre las personas con deficiencias y \textbf{las barreras debidas a la actitud y al entorno}. \emph{Traducido a un ejemplo que se entiende de una:} \textbf{una persona en silla de ruedas no «tiene» la imposibilidad de entrar a la biblioteca ---la imposibilidad la tienen \emph{las escaleras}---.} \emph{Y fíjate en que la Convención nombra dos entornos:} \textbf{el físico y \emph{la actitud}}. \emph{La actitud también es una barrera:} \textbf{hablarle a alguien más despacio de lo necesario, decidir por él, felicitarlo por cosas ordinarias, hablarle al acompañante en vez de a él.}}
\milcpaso{3}{milcTurquesa}{\textbf{Una lengua tardó hasta 1996 en ser lengua.} Los niños sordos colombianos se comunicaban entre sí desde siempre; \textbf{el Estado reconoció la \emph{Lengua Manual Colombiana} como idioma propio de la comunidad sorda con la Ley 324 de 1996}, que además ordenó garantizar intérpretes y promover escuelas para formarlos. \emph{Piensa lo que significa que algo tan evidente necesitara una ley:} \textbf{durante décadas, lo que esos niños hablaban se consideró un apaño, no una lengua} ---y a menudo se les enseñaba a hablar en vez de enseñarles en la suya---. \emph{La lección general vale para cualquier diferencia:} \textbf{muchas veces la barrera no es que alguien no pueda comunicarse ---es que lo suyo no se cuenta como comunicación válida---.}}
\milcpaso{4}{milcTurquesa}{\textbf{Los ajustes que no cuestan plata.} \emph{Es fácil salir de aquí pensando que esto se arregla con obras, y no es así:} \textbf{la mayoría de las barreras de un salón se quitan gratis}. \emph{Ejemplos reales:} \textbf{hablar de frente} para quien lee los labios; \textbf{no explicar mientras se escribe} de espaldas; \textbf{dejar la instrucción por escrito} además de decirla; \textbf{dar un minuto más} para copiar; \textbf{dejar sentarse adelante} a quien lo necesita; \textbf{decir la palabra difícil deletreada}; \textbf{avisar los cambios con anticipación} ---esto último cambia el día entero de muchas personas neurodivergentes---; \textbf{no tocar la silla, el bastón ni al perro guía sin permiso}. \emph{Ninguna cuesta un peso, y casi todas mejoran la clase para todo el mundo.}}

\begin{drawbox}{milcTurquesa}{Las tres clases de barrera}
\begin{pasoslista}
\item \textbf{Física:} escaleras sin rampa, puertas estrechas, baños, tableros que no se ven, ruido.
\item \textbf{De comunicación:} todo dicho y nada escrito, avisos solo por audio, explicar de espaldas, letra pequeña.
\item \textbf{De método:} una sola forma de evaluar, instrucciones largas y habladas, cambios sin avisar, ritmo único.
\item \textbf{Y la que atraviesa las tres ---la actitud---:} decidir por alguien, hablarle al acompañante, felicitarlo por lo ordinario.
\item \textbf{Regla de oro:} \emph{la pregunta no es «¿qué le pasa a esa persona?», sino «¿qué hay aquí que se lo impide?».}
\end{pasoslista}
\end{drawbox}

\begin{softbox}{milcMostaza}{milcGris}{Errores comunes y su corrección}
\textbf{Mirar a la persona y no al lugar:} la Convención dice que la barrera está en el entorno. \textbf{Creer que solo son rampas:} las de comunicación y método son más frecuentes y más invisibles. \textbf{Ayudar sin preguntar:} se pregunta primero, siempre. \textbf{Hablarle al acompañante:} la persona está ahí. \textbf{Felicitar por lo ordinario:} venir a clase no es una hazaña. \textbf{«Actividades especiales aparte»:} aparte no es incluido. \textbf{Suponer que la diferencia se ve:} muchas no se ven.
\end{softbox}

\begin{infoband}{milcTurquesa}


\textbf{En tu cuaderno (Actividad 2 · EXPLICA).} Título: «Actividad 2. Dónde está la barrera». \textbf{Formato:} toma \textbf{cinco} frases del tipo «no puede\ldots» y reescríbelas como «aquí no hay\ldots» o «aquí se hace de tal forma que\ldots». \textbf{Extensión:} media página. \textbf{Sabes que terminaste cuando} puedes explicar por qué \textbf{las escaleras son la barrera y no la silla de ruedas}.
\end{infoband}

\puente{Ya sabes dónde mirar. Es hora de recorrer tu colegio con esos ojos ---y de proponer arreglos que no cuestan plata---.}

\milcestacion{milcMagenta}{Fase 3 · Praxis: auditar el colegio}

\begin{softbox}{milcMagenta}{milcGris}{Cómo se audita un lugar buscando barreras}
Esta guía termina en un papel entregado. Auditar sin entregar es hacer la mitad.
\end{softbox}

\milcpaso{1}{milcMagenta}{\textbf{El recorrido (15 min, en grupos, por el colegio).} Recorran una ruta real ---entrada, salón, baño, biblioteca, cafetería, cancha, la sala de sistemas--- \textbf{anotando barreras}. \emph{Un truco que funciona:} \textbf{háganlo imaginando a cuatro personas distintas} ---alguien que no ve bien, alguien que no oye, alguien que no puede subir escaleras, alguien a quien el ruido y los cambios lo desbordan--- y anoten qué le pasaría a cada una en cada punto.}
\milcpaso{2}{milcMagenta}{\textbf{Las tres clases (8 min).} Clasifiquen lo encontrado en \textbf{física}, \textbf{de comunicación} y \textbf{de método}. \emph{Cuenten cuántas hay de cada una.} \textbf{Casi siempre las físicas son las menos y las de método las más} ---y son, justamente, las que nadie había mirado nunca---.}
\milcpaso{3}{milcMagenta}{\textbf{El ajuste sin plata (10 min).} Al frente de cada barrera, escriban \textbf{un ajuste que no cueste dinero} y \textbf{quién tendría que hacerlo} ---un profesor, la coordinación, los mismos estudiantes---. \emph{Descarten por ahora las que necesitan obra:} \textbf{déjenlas anotadas aparte, pero el producto son las gratis}, porque son las que pueden ocurrir este mes.}
\milcpaso{4}{milcMagenta}{\textbf{Entregar la lista (fuera de clase).} \textbf{Pásenla en limpio y entréguensela a alguien que pueda decidir}: el director de grupo, coordinación, el consejo estudiantil. \emph{Y háganlo bien:} \textbf{sin acusar a nadie} ---las barreras no las puso nadie a propósito, y presentarlas como reclamo garantiza que se defiendan---. \emph{Preséntenlas como lo que son:} \textbf{una lista de cosas fáciles que mejoran el colegio para todos}. \textbf{Anoten qué respondieron.}}

\begin{drawbox}{milcMagenta}{Antes de cerrar tu auditoría}
\begin{checklista}
\item Recorrieron una ruta real y no la imaginaron desde el salón.
\item Hicieron el recorrido pensando en \textbf{cuatro} personas distintas.
\item Tienen barreras de los \textbf{tres tipos}, no solo físicas.
\item Cada barrera tiene un \textbf{ajuste sin plata} y \textbf{quién} tendría que hacerlo.
\item Todas sus frases hablan del lugar, no de las personas.
\item \textbf{Entregaron la lista} y anotaron qué respondieron.
\end{checklista}
\end{drawbox}

\begin{softbox}{milcMagenta}{milcGris}{Plantilla de trabajo}
\textbf{Cómo se escribe una barrera.} \emph{«En \_\_\_\_, quien \_\_\_\_ no puede \_\_\_\_ porque aquí \_\_\_\_.»} ---el sujeto de la última parte es el lugar---. \\[1mm]
\textbf{El ajuste.} \emph{«Se arregla si \_\_\_\_. Lo tendría que hacer \_\_\_\_. No cuesta nada.»} \\[1mm]
\textbf{Al entregar.} \emph{«Hicimos una lista de cosas fáciles que harían el colegio más fácil para todos. ¿Se la podemos mostrar?»}
\end{softbox}

\begin{infoband}{milcMagenta}


\textbf{En tu cuaderno (Actividad 3 · CREA).} Título: «Actividad 3. Auditoría de barreras». \textbf{Formato:} la ruta recorrida + las barreras clasificadas en los tres tipos con su conteo + el ajuste y el responsable de cada una + a quién se entregó y qué respondió. \textbf{Extensión:} una página. \textbf{Sabes que terminaste cuando} la lista \textbf{está entregada}.
\end{infoband}

\puente{El producto se entrega. \emph{Una auditoría que se queda en el cuaderno no le quitó una barrera a nadie.}}

\milcestacion{milcMostaza}{Producto de la sesión}
\textbf{Auditoría de barreras: recorrido real, barreras de los tres tipos con ajustes sin costo, y la lista entregada}.
\begin{softbox}{milcMostaza}{milcGris}{Criterios de calidad}
(1) La auditoría surge de un recorrido real y contempla al menos cuatro perfiles distintos. (2) Las barreras están formuladas como propiedades del entorno, no de las personas. (3) Hay barreras de los tres tipos, incluidas las de comunicación y método. (4) Cada barrera tiene un ajuste sin costo y un responsable identificado. (5) La lista fue entregada a alguien con capacidad de decidir, sin formularla como reclamo.
\end{softbox}

\puente{Antes de cerrar, revisa lo que suele faltar: si tu lista solo tiene escaleras y rampas, te faltaron las barreras de comunicación y de método.}

\milcestacion{milcVino}{Evaluación liberadora · cinco dimensiones}

\begin{softbox}{milcVino}{milcGris}{Sentido de la evaluación}
Evaluar no es solo revisar una nota. Es reconocer qué aprendí, qué cambió en mí, qué emociones manejé, cómo me relaciono con la ciudad y la comunidad, y qué le dejaré a quien viene detrás.
\end{softbox}

\textbf{Mi reflexión liberadora en cinco dimensiones.} Completa cada frase iniciada:

\small
\begin{tabularx}{\textwidth}{>{\bfseries\RaggedRight\arraybackslash}p{3.4cm}Y>{\RaggedRight\arraybackslash}p{4.6cm}}
\rowcolor{milcVino}\color{white}Dimensión & \color{white}\bfseries Mi respuesta (frame starter) & \color{white}\bfseries Algo que descubrí\\
1. Personal & Lo que más cambió en mí fue \linead{6.5cm} & \linead{4.2cm}\\
2. Emocional & La emoción más fuerte fue \linead{2.5cm} y la regulé \linead{3cm} & \linead{4.2cm}\\
3. Ciudadana & En democracia esto importa porque \linead{6cm} & \linead{4.2cm}\\
4. Local (Cartago) & En mi barrio sería útil para \linead{6.5cm} & \linead{4.2cm}\\
5. Intergeneracional & A mi abuela/abuelo le diría \linead{6.5cm} & \linead{4.2cm}\\
\end{tabularx}
\normalsize

\begin{infoband}{milcVino}
\textbf{Para la próxima sustentación.} Vuelve a esta tabla en seis meses. Verás que las respuestas cambian — no porque lo aprendido cambie, sino porque tú cambias. La evaluación liberadora es una conversación contigo a través del tiempo.
\end{infoband}


\puente{Cerramos con tres voces: el rostro que reclama un derecho, la abeja y el enjambre, y el cuidado de un entorno compartido.}

\milcestacion{milcGradeDark}{Triángulo de pensamiento · cierre filosófico}

\begin{softbox}{milcVino}{milcGris}{Dussel · lente del nosotros}
\textit{«El otro se revela realmente como otro\ldots como el pobre, el oprimido; el que, a la vera del camino, fuera del sistema, muestra su rostro sufriente y sin embargo desafiante: «¡Tengo hambre!, ¡tengo derecho a comer!»»}

{\footnotesize\itshape\color{milcNegro!70}--- Enrique Dussel · Filosofía de la liberación (1977)}

\textbf{«Fuera del sistema»} es, en esta guía, una descripción literal: \emph{un edificio con escaleras deja a alguien fuera del sistema sin que nadie lo haya decidido}. \textbf{Y la última palabra vuelve a ser la clave:} \emph{el que reclama una rampa, un intérprete o una instrucción escrita \textbf{no está pidiendo un favor} ---está reclamando algo que Colombia firmó en 2006 y aprobó por ley en 2009---.} \textbf{Ese cambio de marco lo cambia todo:} \emph{una barrera no es una desgracia que le tocó a alguien ---es un incumplimiento---.} \emph{Y por eso la lista que ustedes entregan no es una petición de buena voluntad: es un recordatorio.}

\textbf{¿He pensado alguna vez que hacer accesible algo era «un detalle bonito» y no una obligación?}
\end{softbox}
\linead{\linewidth}\\[1mm]
\linead{\linewidth}

\begin{softbox}{milcMostaza}{milcGris}{Estoicismo · Marco Aurelio · Meditaciones VI, 54 (c. 175 d.C.) · cuidado interior}
\textit{«Lo que no beneficia al enjambre, tampoco beneficia a la abeja.»}

{\footnotesize\itshape\color{milcNegro!70}--- Marco Aurelio · Meditaciones VI, 54 (c. 175 d.C.)}

\emph{Aquí la frase tiene una vuelta que sorprende y que conviene mirar bien:} \textbf{los ajustes que quitan barreras casi siempre mejoran la clase \emph{para todo el mundo}}. \emph{Piénsalo con los ejemplos:} \textbf{dejar la instrucción por escrito ayuda al que no oye \emph{y} al que se distrajo; hablar de frente ayuda al que lee los labios \emph{y} al del último puesto; avisar los cambios con tiempo ayuda a quien los cambios lo desbordan \emph{y} a quien tiene que organizar su tarde.} \textbf{Ese es el mejor argumento de esta guía y no es un argumento moral, es práctico:} \emph{lo que se hace por uno termina sirviéndole al enjambre entero.}

\textbf{De los ajustes de nuestra lista, ¿cuáles me servirían también a mí?}
\end{softbox}
\linead{\linewidth}\\[1mm]
\linead{\linewidth}

\begin{softbox}{milcTurquesa}{milcGris}{Floridi · lente de la infoesfera}
\textit{«El florecimiento de las entidades informacionales y de toda la infoesfera debe promoverse preservando, cultivando y enriqueciendo sus propiedades.»}

{\footnotesize\itshape\color{milcNegro!70}--- Luciano Floridi · La ética de la información (2013; trad.)}

\textbf{Hoy una parte enorme del colegio ocurre en pantallas, y ahí se levantan barreras nuevas sin que nadie lo note.} \emph{Un aviso importante puesto solo como \textbf{imagen} no lo puede leer un lector de pantalla; un video sin subtítulos deja por fuera a quien no oye ---y también a quien no tiene audífonos---; un PDF que es una foto de un texto no se puede buscar ni ampliar.} \textbf{Y la buena noticia es la misma de siempre:} \emph{arreglarlo es gratis} ---escribir el texto además de la imagen, poner subtítulos, describir las fotos---. \textbf{Cultivar, aquí, significa que la información del curso llegue completa a todos los del curso.} \emph{Empezando por el grupo de chat, que también es una puerta.}

\textbf{Los avisos importantes de mi curso, ¿llegan de una forma que todos puedan recibir?}
\end{softbox}
\linead{\linewidth}\\[1mm]
\linead{\linewidth}

\begin{infoband}{milcNegro}
\textbf{Nota para el docente.} Las tres son citas textuales y llevan su referencia debajo. Puedes remitir a la obra en clase.
\end{infoband}

\milcestacion{milcVino}{Mi autoevaluación · marca una casilla por fila}

{\small
\setlength{\extrarowheight}{2.2pt}
\begin{tabularx}{\textwidth}{>{\RaggedRight\arraybackslash\bfseries}p{3.0cm}YYY}
\rowcolor{milcVino}\color{white}\bfseries Criterio & \color{white}\bfseries Lo logré & \color{white}\bfseries En proceso & \color{white}\bfseries Necesito apoyo\\[.6mm]
Comprensión del tema & \milcopcion{Explico las ideas con mis palabras.} & \milcopcion{Reconozco las ideas, falta precisión.} & \milcopcion{Necesito repasar lo básico.}\\[.8mm]
\rowcolor{milcGris}Calidad de la auditoría & \milcopcion{Encontré barreras de los tres tipos y propuse ajustes concretos que no cuestan plata.} & \milcopcion{Identifico rampas y escaleras, pero no las barreras de comunicación ni de método.} & \milcopcion{Sigo pensando que el problema lo tiene la persona y no el lugar.}\\[.8mm]
Aplicación y producto & \milcopcion{Mi producto cumple los criterios.} & \milcopcion{Está armado, con ajustes pendientes.} & \milcopcion{Aún está en construcción.}\\[.8mm]
\rowcolor{milcGris}Reflexión 5 dimensiones & \milcopcion{Respondí cada una con honestidad.} & \milcopcion{Respondí algunas con detalle.} & \milcopcion{Mis respuestas son muy generales.}\\[.8mm]
Triángulo de pensamiento & \milcopcion{Conecté las 3 voces con mi trabajo.} & \milcopcion{Conecté al menos una.} & \milcopcion{No las relaciono con mi proceso.}\\[.8mm]
\end{tabularx}}

\vspace{3mm}
\textbf{Hoy entendí que:}\\[1mm]
\linead{\linewidth}\\[1mm]
\linead{\linewidth}

\textbf{Mi compromiso para la próxima sesión es:}\\[1mm]
\linead{\linewidth}\\[1mm]
\linead{\linewidth}

\begin{softbox}{milcMostaza}{milcGris}{Cita de despedida · Marco Aurelio}
\textit{``El obstáculo en el camino se vuelve el camino. Lo que estorba al camino se vuelve camino.''}\\[1mm]
Cada error de hoy, bien observado, se convierte en pieza del camino que pisarás mañana.
\end{softbox}

\small
\begin{tabularx}{\textwidth}{>{\bfseries}p{3.6cm}Y}
\rowcolor{milcGradeDark!12}Nombre completo: & \linead{10cm}\\
Fecha: & \linead{4cm} \quad Curso/grupo: \linead{4cm}\\
Mi firma: & \linead{9cm}\\
\end{tabularx}
\normalsize

\begin{infoband}{milcGradeDark}
\textbf{Para la próxima sesión.} Dos cosas. \textbf{Primero, entreguen la lista} y traigan \textbf{qué respondieron}. \emph{Y cada uno, por su cuenta, aplique \textbf{un} ajuste de la lista durante la semana ---hablar de frente, escribir lo que se dice, avisar los cambios--- y anote si alguien lo notó.} \textbf{Segundo, pregunta en tu casa qué pasaba antes:} averigua con alguien mayor \textbf{si en su colegio había algún niño que no viera, no oyera, no caminara o aprendiera distinto}, \textbf{qué se hacía con él} ---si iba a clase, si lo mandaban a otro lado, si no iba a estudiar--- y \textbf{qué se decía de esa persona}. \textbf{Trae:} la respuesta a la lista y lo que te contaron. \emph{Prepárate para una respuesta que se repite mucho y que dice más que cualquier estadística: «no había ninguno». Casi nunca era verdad que no existieran. Es que no estaban ahí.}
\end{infoband}

\end{document}
